\chapter{Introduction}
\label{chap:intro}

The term structure of interest rates is one of the most studied objects in quantitative finance \textcolor{red}{how do we know that?}. It encodes the price of money at every maturity, underlies the valuation of bonds, swaps and interest-rate derivatives, and serves as an input to monetary policy, risk management and corporate finance. Two features make it difficult to model. First, the curve is intrinsically high-dimensional --- at any point in time it consists of yields across a continuum of maturities, observed only at finitely many discrete points. Second, any model of the curve must be internally consistent in the sense of absence of arbitrage. These two requirements --- compression of a high-dimensional object and consistency with arbitrage-free pricing --- are non-trivial to reconcile.

Approaches to this problem range from classical parametric specifications to recent neural-network architectures that learn a latent representation directly from data. Chapter~\ref{ch:model-background} reviews this line of work in detail. The starting point for the present thesis is an affine-inspired encoder--decoder model that imposes the no-arbitrage condition as a hard constraint at the decoding stage. The decoder maps a low-dimensional latent state into zero-coupon bond prices through a Markovian factor-pricing PDE, so that the swap rates produced by the model are arbitrage-free by construction.

Building on this architecture, the thesis asks three related questions. First, can the model represent the par swap-rate curve accurately across nine major currencies and a wide range of market regimes, and how robust is its performance out of sample. Second, can the learned latent dynamics be used not only for cross-sectional reconstruction but also for simulation of forward scenarios. Third, can the same dynamics be used to price interest-rate derivatives, specifically European swaptions. The three questions correspond to three distinct empirical objectives --- reconstruction, simulation and pricing --- which are tested in turn. 

The thesis makes three main contributions.

First, the in-sample and rolling out-of-sample evaluation establishes that the baseline arbitrage-free encoder--decoder reconstructs the swap-rate curve well in benign periods but exhibits a flexibility--robustness trade-off at higher latent dimensions. The two-factor model achieves an average in-sample RMSE of approximately twelve basis points across the nine currencies. Increasing the latent dimension to four reduces the in-sample error further but produces rolling out-of-sample errors that are materially worse, particularly during the post-2014 negative-rate era and the post-2022 tightening cycle. The qualitative reconstruction of less common curve shapes --- negative-rate and inverted curves --- is unreliable across all latent dimensions, and augmenting the encoder with simple slope and curvature features improves in-sample fit but does not improve out-of-sample performance.

Second, we propose a \emph{stable} variant of the model in which the latent dynamics are restricted to be mean-reverting, with bounded volatilities, valid correlation matrices and a bounded short rate. Each restriction is financially motivated and individually justifiable. Their joint effect is to render the latent stochastic differential equation well-posed in the sense of Proposition~\ref{prop:sde_existence}, while leaving in-sample reconstruction accuracy at least as good as the baseline. The four-factor stable model becomes the best overall performer, achieving both the lowest in-sample RMSE and the lowest rolling out-of-sample RMSE among all variants tested. This reverses the baseline's apparent capacity--robustness trade-off and unlocks the ability to simulate the latent process forward over multi-year horizons without producing explosive trajectories.

Third, we use the stable model to price European swaptions and find that the unmodified model, although arbitrage-free at any point in time, gives poor option prices when simulated forward, with an average ATM straddle implied-vol error of around 415 basis points. A minimal recalibration layer with eight parameters --- the \emph{Constant Market Price of Risk} layer --- reduces this to approximately 42 basis points out of sample, demonstrating that the latent state contains useful information for pricing once the risk-neutral dynamics are recalibrated. More elaborate yield-curve-only pricing layers overfit and worsen out-of-sample performance, suggesting that the yield curve alone does not identify the option-implied volatility regime. When the pricing layer is augmented with leave-one-cell-out features extracted from the observed swaption volatility surface, the test error falls further to approximately 32 basis points. The roughly 10 bp gap between the two models can be interpreted as the value of contemporaneous option-market volatility information that is not contained in the yield curve.

The remainder of the thesis is organised as follows. Chapters~\ref{chap:finance_theory} and~\ref{chap:ml_theory} develop the financial and machine-learning theory required for what follows. Chapter~\ref{ch:model-background} places the affine-inspired encoder--decoder architecture in the term-structure modelling literature, and Chapter~\ref{chap:encoderdecodermodel} describes the architecture in detail. Chapter~\ref{chap:datadescription} introduces the Bloomberg swap-rate dataset used throughout. Chapter~\ref{chap:evaluation} evaluates reconstruction accuracy and out-of-sample robustness for latent dimensions \(\ell\in\{1,2,3,4\}\). Chapter~\ref{chap:simulation} introduces the stable variant and compares the simulated latent paths of the two model specifications. Chapter~\ref{chap:pricing_methodology} treats swaption pricing, including the Constant MPR and Surface Vol MPR layers introduced above and a brief discussion of the model-implied volatility smile. Chapter~\ref{ch:discussion} synthesises the findings, and Chapter~\ref{chap:con} concludes.
